\documentclass[conference,compsoc]{IEEEtran}
% Some/most Computer Society conferences require the compsoc mode option,
% but others may want the standard conference format.
%
% If IEEEtran.cls has not been installed into the LaTeX system files,
% manually specify the path to it like:
% \documentclass[conference,compsoc]{../sty/IEEEtran}

\usepackage{tcse}


\begin{document}
%
% paper title
% Titles are generally capitalized except for words such as a, an, and, as,
% at, but, by, for, in, nor, of, on, or, the, to and up, which are usually
% not capitalized unless they are the first or last word of the title.
% Linebreaks \\ can be used within to get better formatting as desired.
% Do not put math or special symbols in the title.
\title{2023台灣軟體工程研討會論文格式\\
Camera Ready Format for the \\
2023 Taiwan Conference on Software Engineering
\thanks{本研究接受國科會編號：NSTC XXX-XXXX-X-XXX-XXX研究計畫經費補助}
}

% author names and affiliations
\author{\IEEEauthorblockN{作者A\textsuperscript{\footnotesize 1}
\hspace{.5em}
作者B\textsuperscript{\footnotesize 2}
\hspace{.5em}
作者C\textsuperscript{\footnotesize 1}\IEEEauthorrefmark{1}}
\IEEEauthorblockA{\textsuperscript{\upshape\footnotesize 1} 服務學校或機關及其部門}
\IEEEauthorblockA{\textsuperscript{\upshape\footnotesize 2} 服務學校或機關及其部門}
\IEEEauthorblockA{\upshape Authors and Affiliation}
\IEEEauthorrefmark{1}Email: name@univ.edu.tw}


% use a multiple column layout for up to three different
% affiliations
%\author{\IEEEauthorblockN{Michael Shell}
%\IEEEauthorblockA{School of Electrical and\\Computer Engineering\\
%Georgia Institute of Technology\\
%Atlanta, Georgia 30332--0250\\
%Email: http://www.michaelshell.org/contact.html}
%\and
%\IEEEauthorblockN{Homer Simpson}
%\IEEEauthorblockA{Twentieth Century Fox\\
%Springfield, USA\\
%Email: homer@thesimpsons.com}
%\and
%\IEEEauthorblockN{James Kirk\\ and Montgomery Scott}
%\IEEEauthorblockA{Starfleet Academy\\
%San Francisco, California 96678-2391\\
%Telephone: (800) 555--1212\\
%Fax: (888) 555--1212}}

% conference papers do not typically use \thanks and this command
% is locked out in conference mode. If really needed, such as for
% the acknowledgment of grants, issue a \IEEEoverridecommandlockouts
% after \documentclass

% for over three affiliations, or if they all won't fit within the width
% of the page (and note that there is less available width in this regard for
% compsoc conferences compared to traditional conferences), use this
% alternative format:
% 
%\author{\IEEEauthorblockN{Michael Shell\IEEEauthorrefmark{1},
%Homer Simpson\IEEEauthorrefmark{2},
%James Kirk\IEEEauthorrefmark{3}, 
%Montgomery Scott\IEEEauthorrefmark{3} and
%Eldon Tyrell\IEEEauthorrefmark{4}}
%\IEEEauthorblockA{\IEEEauthorrefmark{1}School of Electrical and Computer Engineering\\
%Georgia Institute of Technology,
%Atlanta, Georgia 30332--0250\\ Email: see http://www.michaelshell.org/contact.html}
%\IEEEauthorblockA{\IEEEauthorrefmark{2}Twentieth Century Fox, Springfield, USA\\
%Email: homer@thesimpsons.com}
%\IEEEauthorblockA{\IEEEauthorrefmark{3}Starfleet Academy, San Francisco, California 96678-2391\\
%Telephone: (800) 555--1212, Fax: (888) 555--1212}
%\IEEEauthorblockA{\IEEEauthorrefmark{4}Tyrell Inc., 123 Replicant Street, Los Angeles, California 90210--4321}}




% use for special paper notices
%\IEEEspecialpapernotice{(Invited Paper)}




% make the title area
\maketitle

% As a general rule, do not put math, special symbols or citations
% in the abstract
\begin{abstract}
本頁內容為有關第十九屆台灣軟體工程研討會的投稿與完稿格式，請投稿人直接依據本格式製作稿件。本文件即依此格式製作完成。
\end{abstract}

% keywords
\begin{IEEEkeywords}
關鍵字1、關鍵字2、關鍵字3
\end{IEEEkeywords}



% For peer review papers, you can put extra information on the cover
% page as needed:
% \ifCLASSOPTIONpeerreview
% \begin{center} \bfseries EDICS Category: 3-BBND \end{center}
% \fi
%
% For peerreview papers, this IEEEtran command inserts a page break and
% creates the second title. It will be ignored for other modes.
\IEEEpeerreviewmaketitle



\section{論文規定}
軟體工程\cite{b1,b2}研討會論文之使用語言為中文或英文，使用紙張大小為A4紙張，不列印頁碼。一般論文(包含中英文)以4-8頁為限，軟體展示論文與海報論文則以2-8頁為限。
\medskip

本稿為台灣軟體工程研討會投稿論文之 \LaTeX 格式範本，基於
IEEE Demo Template for Computer Society Conferences
進行修改而成。要求IEEEtran.cls 1.8b之後的版本，並須以XeLaTeX進行編譯，
建議使用Overleaf線上服務進行排版。


\section{論文格式}
主要論文格式如下：
\begin{itemize}
\item 單行間距(single space)
\item 雙欄文章(two column format)
\item 第一頁開頭依序為中文標題、英文標題、中文作者姓名與服務單位，及英文作者姓名與服務單位名稱，摘要與論文內容。
\end{itemize}

\section{圖表與公式}
圖表內之文字請與內容相符，圖下或表上請以連續號碼編號，並加註圖表號碼，如: 圖1或表1。公式請說明符號之意義，並請加上連續標號。圖表與公式之範例如下:

\begin{figure}[htbp]
\centerline{\includegraphics{fig1.png}}
\caption{圖範例}
\label{fig1}
\end{figure}


\begin{table}[htbp]
\caption{表範例}
\begin{center}
\begin{tabular}{ccccc}
\hline
\textbf{Mar 2016} & \textbf{Mar 2015} & \textbf{Change} & \textbf{Language} & \textbf{Ratings} \\
\hline
1 & 2 & $\Uparrow$   & Java   & 20.53\% \\
\hline
2 & 1 & $\Uparrow$   & C      & 14.60\% \\
\hline
3 & 4 & $\Downarrow$ & C++    & 6.72\%  \\
\hline
4 & 5 & $\Uparrow$   & C\#    & 4.27\%  \\
\hline
5 & 8 & $\Uparrow$   & Python & 4.26\%  \\
\hline
6 & 6 &              & PHP    & 2.77\%  \\
\hline
\end{tabular}
\label{tab1}
\end{center}
\end{table}

\begin{equation}
UFC = \sum_{i=1}^{5} \sum_{j=1}^{3} N_{ij} W_{ij} \label{eq1}
\end{equation}

其中，$N_{ij}$ 與 $W_{ij}$ 分別代表分類$i$與複雜度$j$之數量與權重。
\medskip

於文中提及上述圖表與公式時，可以圖\ref{fig1}、表\ref{tab1}、公式\eqref{eq1}的形式參照。


\section{其它}
其它未詳細規定之部份請以一般國際會議論文格式製作。



% An example of a floating figure using the graphicx package.
% Note that \label must occur AFTER (or within) \caption.
% For figures, \caption should occur after the \includegraphics.
% Note that IEEEtran v1.7 and later has special internal code that
% is designed to preserve the operation of \label within \caption
% even when the captionsoff option is in effect. However, because
% of issues like this, it may be the safest practice to put all your
% \label just after \caption rather than within \caption{}.
%
% Reminder: the "draftcls" or "draftclsnofoot", not "draft", class
% option should be used if it is desired that the figures are to be
% displayed while in draft mode.
%
%\begin{figure}[!t]
%\centering
%\includegraphics[width=2.5in]{myfigure}
% where an .eps filename suffix will be assumed under latex, 
% and a .pdf suffix will be assumed for pdflatex; or what has been declared
% via \DeclareGraphicsExtensions.
%\caption{Simulation results for the network.}
%\label{fig_sim}
%\end{figure}

% Note that the IEEE typically puts floats only at the top, even when this
% results in a large percentage of a column being occupied by floats.


% An example of a double column floating figure using two subfigures.
% (The subfig.sty package must be loaded for this to work.)
% The subfigure \label commands are set within each subfloat command,
% and the \label for the overall figure must come after \caption.
% \hfil is used as a separator to get equal spacing.
% Watch out that the combined width of all the subfigures on a 
% line do not exceed the text width or a line break will occur.
%
%\begin{figure*}[!t]
%\centering
%\subfloat[Case I]{\includegraphics[width=2.5in]{box}%
%\label{fig_first_case}}
%\hfil
%\subfloat[Case II]{\includegraphics[width=2.5in]{box}%
%\label{fig_second_case}}
%\caption{Simulation results for the network.}
%\label{fig_sim}
%\end{figure*}
%
% Note that often IEEE papers with subfigures do not employ subfigure
% captions (using the optional argument to \subfloat[]), but instead will
% reference/describe all of them (a), (b), etc., within the main caption.
% Be aware that for subfig.sty to generate the (a), (b), etc., subfigure
% labels, the optional argument to \subfloat must be present. If a
% subcaption is not desired, just leave its contents blank,
% e.g., \subfloat[].


% An example of a floating table. Note that, for IEEE style tables, the
% \caption command should come BEFORE the table and, given that table
% captions serve much like titles, are usually capitalized except for words
% such as a, an, and, as, at, but, by, for, in, nor, of, on, or, the, to
% and up, which are usually not capitalized unless they are the first or
% last word of the caption. Table text will default to \footnotesize as
% the IEEE normally uses this smaller font for tables.
% The \label must come after \caption as always.
%
%\begin{table}[!t]
%% increase table row spacing, adjust to taste
%\renewcommand{\arraystretch}{1.3}
% if using array.sty, it might be a good idea to tweak the value of
% \extrarowheight as needed to properly center the text within the cells
%\caption{An Example of a Table}
%\label{table_example}
%\centering
%% Some packages, such as MDW tools, offer better commands for making tables
%% than the plain LaTeX2e tabular which is used here.
%\begin{tabular}{|c||c|}
%\hline
%One & Two\\
%\hline
%Three & Four\\
%\hline
%\end{tabular}
%\end{table}


% Note that the IEEE does not put floats in the very first column
% - or typically anywhere on the first page for that matter. Also,
% in-text middle ("here") positioning is typically not used, but it
% is allowed and encouraged for Computer Society conferences (but
% not Computer Society journals). Most IEEE journals/conferences use
% top floats exclusively. 
% Note that, LaTeX2e, unlike IEEE journals/conferences, places
% footnotes above bottom floats. This can be corrected via the
% \fnbelowfloat command of the stfloats package.




% conference papers do not normally have an appendix



% use section* for acknowledgment
%\ifCLASSOPTIONcompsoc
%  % The Computer Society usually uses the plural form
%  \section*{Acknowledgments}
%\else
%  % regular IEEE prefers the singular form
%  \section*{Acknowledgment}
%\fi
%
%The authors would like to thank...





% trigger a \newpage just before the given reference
% number - used to balance the columns on the last page
% adjust value as needed - may need to be readjusted if
% the document is modified later
%\IEEEtriggeratref{8}
% The "triggered" command can be changed if desired:
%\IEEEtriggercmd{\enlargethispage{-5in}}

% references section

% can use a bibliography generated by BibTeX as a .bbl file
% BibTeX documentation can be easily obtained at:
% http://mirror.ctan.org/biblio/bibtex/contrib/doc/
% The IEEEtran BibTeX style support page is at:
% http://www.michaelshell.org/tex/ieeetran/bibtex/
%\bibliographystyle{IEEEtran}
% argument is your BibTeX string definitions and bibliography database(s)
%\bibliography{IEEEabrv,../bib/paper}
%
% <OR> manually copy in the resultant .bbl file
% set second argument of \begin to the number of references
% (used to reserve space for the reference number labels box)
\begin{thebibliography}{1}

\bibitem{b1}
W.~S. Humphrey, \emph{A discipline for software engineering:} Addison-Wesley Longman Publishing Co., Inc., 1995.

\bibitem{b2}
B.~Kitchenham, O.~Pearl Brereton, D.~Budgen, M.~Turner, J.~Bailey, and S.~Linkman, "Systematic literature reviews in software engineering – A systematic literature review," \emph{Information and Software Technology,} vol. 51, pp. 7-15, 2009.

%\bibitem{IEEEhowto:kopka}
%H.~Kopka and P.~W. Daly, \emph{A Guide to \LaTeX}, 3rd~ed.\hskip %1em plus
%  0.5em minus 0.4em\relax Harlow, England: Addison-Wesley, 1999.

\end{thebibliography}




% that's all folks
\end{document}


